\section{Preliminaries}
\label{sec:background}

\subsection{Concurrency Control}
\subsection{Timestamp-based Protocols}
\subsubsection{Serial Safety Net}

\subsection{Write Ahead Logging for Recovery}

\subsection{Logging}
\subsubsection{P-WAL}

\begin{algorithm}[t]
\caption{\bfseries Per-Thread WAL Commit for Update Transactions (P-WAL)}
\label{alg:pwal}
\begin{algorithmic}[1]
\State \textbf{Shared:} global LSN counter $L$; per-stream durable LSN
       $f[i]$; global durable prefix $P = \min_i f[i]$
\Function{PwalCommit}{$T$, worker $w$}
    \State $\mathit{records} \gets$ \Call{BuildLogRecords}{$T$}
           \Comment{writes and commit record}
    \ForAll{$r \in \mathit{records}$}
        \State $r.\mathit{lsn} \gets$ \Call{FetchInc}{$L$}
               \Comment{global LSN allocation}
        \State \Call{Append}{$\mathit{stream}[w],\, r$}
               \Comment{worker-private WAL stream}
    \EndFor
    \State $T.\ell \gets \max \{\, r.\mathit{lsn} \mid r \in \mathit{records} \,\}$
           \Comment{commit LSN}
    \State \Call{Write}{$\mathit{stream}[w]$};
           \Call{Fdatasync}{$\mathit{stream}[w]$}
           \Comment{worker waits for persistence}
    \State $f[w] \gets \max(f[w],\, T.\ell)$
    \State $P \gets \min_i f[i]$
           \Comment{global durable prefix}
    \While{$P < T.\ell$}
        \State $P \gets \min_i f[i]$
    \EndWhile
    \State \Call{Acknowledge}{$T$}
\EndFunction
\end{algorithmic}
\end{algorithm}


% Problem of P-WAL
P-WAL is attractive because it removes the single-WAL insertion bottleneck by assigning transactions to multiple log streams, which we hereafter call \emph{WAL shards}.  However, combining P-WAL with SSN exposes three
problems.

\textbf{P1. Multiple Global Counter Accesses:~} First, P-WAL can introduce frequent global-counter accesses.  Even if the CC
layer already assigns a commit timestamp for SSN validation and serialization, the recovery layer may allocate a separate global LSN or maintain a global durable prefix.  This duplicates the ordering work already performed by the CC layer and adds shared-memory synchronization to the commit path.

\textbf{P2. Waiting for Data Transfer:~} Second, workers can spend a long time waiting around \texttt{fdatasync}.  P-WAL parallelizes log streams, but a transaction cannot be acknowledged until the corresponding durable condition is satisfied.  If worker threads wait directly
for log flushes or durable-prefix updates, storage synchronization and notification latency reduce the effective scalability of the CC protocol.

\textbf{P3. Global-prefix over-waiting:~} Third, global-prefix acknowledgment waits for log records that are unrelated to the transaction being acknowledged.  This rule is safe, but overly conservative:
if one WAL shard is slow, transactions on other shards may be delayed even when they have no dependency on transactions logged by the slow shard.  The opposite design point, local-only acknowledgment, avoids this delay but is unsafe.  If transaction $T$ reads a value written by transaction $U$, acknowledging $T$ before $U$ is recoverable may allow a crash recovery state that contains $T$ without the state on which $T$ depends.
